%% kapa: exam=1 round=20 year=2009 period=2 form=A subject=appraisal_law subject_ko=부동산관계법규 seq=3 number=43
%% source: Q-Net 감정평가사 기출문제 (cst00309) artlSeq=5024178 file=extracted/1008/5024178/2009년 제20회 감정평가사 1차시험 기출문제/2교시.hwp
%% conversion: kordoc 4.24.51 -> tools/gen1.py (auto); review=proofread-sonnet; flags=none
\begin{question}{43}
국토의 계획 및 이용에 관한 법령상 도시관리계획의 수립절차에 관한 설명으로 옳지 \uline{않은} 것은?
\begin{choices}
  \item 국토해양부장관은 국가계획과 관련된 경우 도시관리계획을 입안할 수 있다.
  \item 도시관리계획은 시ㆍ도지사가 직접 또는 시장ㆍ군수의 신청에 따라 결정하지만, 개발제한구역․수산자원보호구역의 지정 및 변경에 관한 도시관리계획은 국토해양부장관이 결정한다.
  \item 국토해양부장관, 시ㆍ도지사, 시장 또는 군수는 도시관리계획을 입안하려면 대통령령으로 정하는 사항에 대하여 해당 지방의회의 의견을 들어야 한다.
  \item 주민은 기반시설의 설치ㆍ정비 또는 개량에 관한 사항에 대하여 도시관리계획을 입안할 수 있는 자에게 도시관리계획의 입안을 제안할 수 있다.
  \item 도시관리계획의 입안 제안신청을 반려하는 행위는 항고소송의 대상이 되는 행정처분에 해당한다는 것이 판례의 입장이다.
\end{choices}
\end{question}
\answer{2}
